\documentclass[12pt,a4paper,oneside]{report}

\usepackage{fontspec}
\usepackage{polyglossia}
\setdefaultlanguage{vietnamese}
\setmainfont{TeX Gyre Termes}
\setsansfont{TeX Gyre Heros}
\setmonofont{TeX Gyre Cursor}
\usepackage[a4paper,margin=2.5cm]{geometry}
\usepackage{graphicx}
\usepackage{booktabs,longtable,tabularx,array,float}
\usepackage{xcolor,listings,amsmath,amssymb,enumitem,fancyhdr}
\usepackage{tikz}
\usetikzlibrary{arrows.meta,positioning,shapes.geometric}
\usepackage[most]{tcolorbox}
\usepackage{hyperref}

\hypersetup{colorlinks=true,linkcolor=blue!50!black,urlcolor=blue!60!black}
\pagestyle{fancy}
\fancyhf{}
\fancyhead[L]{IMOVEV2}
\fancyhead[R]{Báo cáo đồ án}
\fancyfoot[C]{\thepage}
\setlength{\headheight}{15pt}
\setlength{\parindent}{1.2cm}
\setlength{\parskip}{0.35em}

\newcommand{\dienvao}{\textbf{Điền vào đây}}
\newcommand{\tacgia}{%
  \begin{tcolorbox}[colback=blue!3,colframe=blue!25,boxrule=0.4pt,arc=1mm]
    \small\textbf{ID và họ tên tác giả/reviewer: Điền vào đây}
  \end{tcolorbox}
}
\newcommand{\danhinh}[1]{%
  \begin{figure}[H]
    \centering
    \begin{tcolorbox}[width=0.88\textwidth,height=5cm,valign=center,colback=gray!5,colframe=gray!45]
      \centering\Large\textbf{Dán hình ảnh vào đây}
    \end{tcolorbox}
    \caption{#1}
  \end{figure}
}
\newcommand{\code}[1]{\texttt{\detokenize{#1}}}

\tikzset{
  nodebox/.style={draw,rounded corners,align=center,minimum width=2.7cm,minimum height=0.9cm,fill=blue!5},
  agent/.style={nodebox,fill=green!8,draw=green!55!black},
  external/.style={nodebox,fill=orange!8,draw=orange!60!black},
  data/.style={nodebox,cylinder,shape border rotate=90,aspect=0.25,fill=purple!7,draw=purple!55!black},
  arrow/.style={-{Latex[length=2mm]},thick}
}

\begin{document}

\begin{titlepage}
  \centering
  {\Large\textbf{Tên trường và khoa: Điền vào đây}\par}
  \vspace{1.5cm}
  {\Large Mã và tên môn học: \dienvao\par}
  \vspace{2cm}
  {\Huge\bfseries IMOVEV2\par}
  \vspace{0.5cm}
  {\Large Hệ thống lập kế hoạch giao thông công cộng thích nghi cho du khách tại Singapore\par}
  \vspace{1.5cm}
  \begin{tabular}{rl}
    Học kỳ: & \dienvao \\
    Mã lớp: & \dienvao \\
    Mã nhóm: & \dienvao \\
    Giảng viên và trợ giảng: & \dienvao \\
  \end{tabular}
  \vfill
  \begin{tabular}{ll}
    ID thành viên 1: \dienvao & Họ tên: \dienvao \\
    ID thành viên 2: \dienvao & Họ tên: \dienvao \\
    ID thành viên 3: \dienvao & Họ tên: \dienvao \\
    ID thành viên 4: \dienvao & Họ tên: \dienvao \\
  \end{tabular}
  \vfill
  {\large Ngày cập nhật gần nhất: 10/06/2026\par}
\end{titlepage}

\pagenumbering{roman}
\tableofcontents
\listoffigures
\listoftables
\clearpage
\pagenumbering{arabic}

\chapter{Group Members}
\section{Thành viên và vai trò}
\tacgia
\begin{longtable}{p{2.5cm}p{3.5cm}p{4cm}p{5cm}}
\toprule
\textbf{Student ID} & \textbf{Full Name} & \textbf{Email} & \textbf{Roles / Responsibilities} \\
\midrule
Điền vào đây & Điền vào đây & Điền vào đây & Backend infrastructure và external APIs; chỉnh sửa nếu cần. \\
Điền vào đây & Điền vào đây & Điền vào đây & Planning, Adaptation và Memory Agent; chỉnh sửa nếu cần. \\
Điền vào đây & Điền vào đây & Điền vào đây & Frontend core flow, authentication và planner; chỉnh sửa nếu cần. \\
Điền vào đây & Điền vào đây & Điền vào đây & Map, realtime alerts và frontend integration; chỉnh sửa nếu cần. \\
\bottomrule
\end{longtable}

\section{Nguyên tắc phân công}
\tacgia
Nhóm phân chia công việc theo ranh giới kỹ thuật: infrastructure và external API, agent/business logic, frontend core, map và realtime. Các thay đổi liên tầng cần review chéo vì contract giữa Pydantic model, FastAPI router và React component liên hệ trực tiếp.

\chapter{Idea}
\section{Problem Statement}
\tacgia
Singapore có mạng lưới MRT, LRT, bus và walking link dày đặc. Du khách lần đầu thường đã có danh sách địa điểm muốn tham quan nhưng chưa biết cách phân bổ chúng vào nhiều ngày, sắp thứ tự hợp lý hoặc chọn phương tiện cân bằng thời gian, chi phí và số lần chuyển tuyến. Các công cụ bản đồ phổ biến giải tốt câu hỏi đi từ A đến B ngay lập tức, nhưng không giải quyết đồng thời bài toán lịch trình nhiều ngày, thích nghi khi điều kiện thay đổi và học từ phản hồi cá nhân.

\section{Product Vision}
\tacgia
IMOVEV2 hướng tới một trợ lý lập kế hoạch giao thông công cộng đồng hành trước và trong chuyến đi. Hệ thống biến danh sách địa điểm thành lịch trình theo ngày, đính kèm route đa phương thức, phản ứng trước mưa hoặc gián đoạn tàu và lưu preference cho người dùng đăng nhập.

\section{Target Users}
\tacgia
Người dùng chính là solo traveler hoặc nhóm nhỏ từ hai đến bốn người, lần đầu hoặc ít kinh nghiệm di chuyển tại Singapore. Họ cần hướng dẫn rõ ràng, muốn biết trước thời gian và chi phí nhưng không muốn tự phân tích toàn bộ mạng lưới giao thông.

\section{Project Goals and Scope}
\tacgia
\begin{itemize}
  \item Tạo itinerary nhiều ngày từ danh sách địa điểm.
  \item Cung cấp route Metro, Bus, Walk, Cycle và Grab khi phù hợp.
  \item Xếp hạng mode theo preference và bối cảnh.
  \item Tạo proposal thích nghi khi có cảnh báo giao thông hoặc thời tiết.
  \item Hoạt động ở guest mode; bổ sung persistence và learning khi đăng nhập.
\end{itemize}
Phiên bản hiện tại tập trung vào Singapore, phụ thuộc OneMap cho routing, không tích hợp thanh toán và chưa hỗ trợ multi-city.

\chapter{Problem Analysis and Decomposition}
\section{Phân rã bài toán}
\tacgia
\begin{longtable}{p{3cm}p{3cm}p{4.5cm}p{3cm}p{3cm}}
\toprule
\textbf{Subproblem} & \textbf{Input} & \textbf{Processing} & \textbf{Output} & \textbf{Module} \\
\midrule
Itinerary setup & Ngày, ngân sách, địa điểm, khách sạn & Kiểm tra form và tạo request & Trip request & React Planner \\
Place validation & Place IDs hoặc tên tự do & Curated lookup, Gemini resolve, startup validation & Place records & Planning Agent \\
Multi-day scheduling & Places và số ngày & Greedy hoặc giữ thứ tự & Day groups & Planning Agent \\
Multimodal routing & Các cặp địa điểm & Gọi OneMap song song & Alternatives & OneMap Service \\
Route ranking & Alternatives, profile, context & Normalize và weighted sum & Recommended mode & Scoring Service \\
Live adaptation & Alert và TripPlan & Reroute/swap và tính delta & Proposal & Adaptation Agent \\
Preference learning & Feedback lặp lại & Rule-based pattern counting & Preference update & Memory Agent \\
Persistence & Trip, alert, feedback & Auth và insert/update & Dữ liệu lưu trữ & Supabase \\
\bottomrule
\end{longtable}

\section{Áp dụng Computational Thinking}
\tacgia
\begin{itemize}
  \item \textbf{Decomposition:} chia hệ thống thành frontend, routers, agents, services, models và persistence.
  \item \textbf{Pattern Recognition:} nhận diện cấu trúc trip/day/leg, nearest-neighbor, context và hành vi đổi mode.
  \item \textbf{Abstraction:} đóng gói route, context và preference thành model/service.
  \item \textbf{Algorithm Design:} xây dựng greedy scheduling, weighted scoring, adaptation và fallback.
\end{itemize}

\chapter{System Overview}
\section{Main Stakeholders and Actors}
\tacgia
Các bên liên quan gồm du khách, nhóm phát triển, nhà cung cấp dữ liệu/dịch vụ và hội đồng đánh giá học thuật. Actor chính gồm du khách lần đầu, người dùng có preference, người đang thực hiện hành trình và người dùng đã đăng nhập.

\section{Core Features}
\tacgia
\begin{enumerate}
  \item Lập lịch trình nhiều ngày và nhiều địa điểm.
  \item Tối ưu thứ tự bằng greedy heuristic có time-window constraints.
  \item So sánh và lựa chọn Metro, Bus, Walk, Cycle và Grab.
  \item Chấm điểm thích nghi theo preference, mưa và giờ cao điểm.
  \item Thích nghi với cảnh báo LTA và thời tiết.
  \item Học preference từ feedback lặp lại.
  \item Gemini hỗ trợ NLP, suggestion, warning và chatbot.
\end{enumerate}

\section{Overall Architecture}
\tacgia
\begin{figure}[H]
\centering
\begin{tikzpicture}[node distance=1.2cm and 1.2cm]
  \node[nodebox] (user) {Traveller};
  \node[nodebox,right=of user] (web) {React + Vite\\Frontend};
  \node[nodebox,right=of web] (api) {FastAPI\\Routers};
  \node[agent,above right=0.5cm and 1.1cm of api] (agents) {Planning / Adaptation\\Memory / Chat};
  \node[agent,below right=0.5cm and 1.1cm of api] (services) {Routing / Transit\\Weather / AI / Scoring};
  \node[data,right=of agents] (db) {Supabase};
  \node[external,right=of services] (ext) {External APIs};
  \draw[arrow] (user) -- (web);
  \draw[arrow] (web) -- (api);
  \draw[arrow] (api) -- (agents);
  \draw[arrow] (api) -- (services);
  \draw[arrow] (agents) -- (db);
  \draw[arrow] (services) -- (ext);
\end{tikzpicture}
\caption{Kiến trúc tổng quan của IMOVEV2.}
\end{figure}

\chapter{Pattern Recognition}
\section{Các pattern được áp dụng}
\tacgia
\begin{longtable}{p{3cm}p{4cm}p{4cm}p{4cm}}
\toprule
\textbf{Pattern} & \textbf{Mẫu nhận diện} & \textbf{Implementation} & \textbf{Ý nghĩa} \\
\midrule
Validation & Place phải có key và thời gian hợp lệ & \code{_REQUIRED_KEYS}, \code{_validate_time()} & Phát hiện data lỗi \\
Structural & Trip gồm Day; Day gồm Leg & \code{TripPlan}, \code{DayPlan}, \code{LegResponse} & Contract nhất quán \\
Scheduling & Điểm gần và phù hợp time window & \code{_day_bucketed_greedy()} & Tạo lịch thực dụng \\
Context & Mưa và giờ cao điểm đổi preference hiệu dụng & \code{ContextSnapshot} & Recommendation thích nghi \\
Ranking & So sánh mode trên bốn dimension & \code{score_alternatives()} & Chọn mode \\
Behavioral & Thay đổi mode lặp lại & \code{learn_from_implicit()} & Học preference \\
Resilience & Mode/API có thể không khả dụng & \code{NoRouteError}, estimates & Graceful degradation \\
\bottomrule
\end{longtable}

\section{Greedy Scheduling Pattern}
\tacgia
Bài toán có nét tương đồng với nearest-neighbor routing, nhưng phải xét số ngày, thời lượng tham quan, opening hours, best-time window và địa điểm buổi tối. Planning Agent dùng heuristic \code{_day_bucketed_greedy()} thay vì exhaustive search.

\begin{figure}[H]
\centering
\begin{tikzpicture}[node distance=0.85cm]
  \node[nodebox] (a) {Classify places};
  \node[nodebox,below=of a] (b) {Create day buckets};
  \node[nodebox,below=of b] (c) {Filter feasible candidates};
  \node[nodebox,below=of c] (d) {Pick nearest candidate};
  \node[nodebox,below=of d] (e) {Update clock and position};
  \node[nodebox,below=of e] (f) {Assign evening places};
  \draw[arrow] (a) -- (b);
  \draw[arrow] (b) -- (c);
  \draw[arrow] (c) -- (d);
  \draw[arrow] (d) -- (e);
  \draw[arrow] (e.east) -- ++(1,0) |- (c.east);
  \draw[arrow] (e) -- (f);
\end{tikzpicture}
\caption{Pattern lập lịch tham lam có ràng buộc.}
\end{figure}

\section{Context-aware Ranking and Behavioral Learning}
\tacgia
Mỗi alternative được mô tả bằng duration, cost, walking minutes và transfers. Các dimension được chuẩn hóa tương đối, sau đó tính weighted sum. Memory Agent nhận diện ít nhất hai lần đổi BUS sang MRT hoặc đổi sang WALK để cập nhật preference có cấu trúc.

\chapter{Abstraction}
\section{Các tầng abstraction}
\tacgia
\begin{table}[H]
\centering
\begin{tabularx}{\textwidth}{p{3.5cm}X X}
\toprule
\textbf{Layer} & \textbf{Implementation} & \textbf{Chi tiết được che giấu} \\
\midrule
Frontend & Pages, components, hooks, API service & HTTP, auth header, cache \\
Router & trips, places, alerts, transit, preferences, chat & HTTP request/response \\
Agent & Planning, Adaptation, Memory, Chat & Use-case orchestration \\
Service & OneMap, LTA, OpenWeather, Gemini, Scoring & API protocol và thuật toán \\
Model & Pydantic models & Dictionary shape không kiểm soát \\
Persistence & Supabase, in-memory, localStorage & Chi tiết lưu trữ \\
\bottomrule
\end{tabularx}
\caption{Các tầng abstraction trong IMOVEV2.}
\end{table}

\section{Hybrid AI-assisted Planning}
\tacgia
Code xác định chịu trách nhiệm validation, scheduling, routing, scoring và fallback. Gemini hỗ trợ resolve tên địa điểm, suggestion, warning và chatbot. Ranh giới này hạn chế hallucination ảnh hưởng route, chi phí hoặc thời gian.

\section{Giới hạn của Abstraction hiện tại}
\tacgia
Hệ thống chưa có interface chung cho nhiều AI provider, chưa phải distributed microservice architecture, một phần persistence vẫn nằm trong routers và curated dataset được Planning Agent load trực tiếp.

\chapter{System / Algorithm Design}
\section{Technical Stack}
\tacgia
\begin{longtable}{p{3.6cm}p{10.5cm}}
\toprule
\textbf{Layer} & \textbf{Technology} \\
\midrule
Frontend & React 18, Vite, React Router, Tailwind CSS, Radix UI, React Leaflet \\
Backend & Python, FastAPI, Pydantic, APScheduler \\
Database/Auth & Supabase PostgreSQL, Auth, RLS, Realtime \\
Routing & OneMap geocoding và multimodal routing \\
Transit & LTA DataMall \\
Weather & OpenWeather \\
AI & Google Gemini 2.5 Flash \\
Testing & Pytest, Vitest, Testing Library \\
\bottomrule
\end{longtable}

\section{Planning Agent Algorithm}
\tacgia
\begin{figure}[H]
\centering
\begin{tikzpicture}[node distance=0.8cm]
  \node[nodebox] (a) {Trip request};
  \node[nodebox,below=of a] (b) {Validate/resolve places};
  \node[nodebox,below=of b] (c) {Schedule across days};
  \node[nodebox,below=of c] (d) {Fetch alternatives};
  \node[nodebox,below=of d] (e) {Normalize and score};
  \node[nodebox,below=of e] (f) {Apply fallback guards};
  \node[nodebox,below=of f] (g) {Build TripPlan};
  \draw[arrow] (a) -- (b);
  \draw[arrow] (b) -- (c);
  \draw[arrow] (c) -- (d);
  \draw[arrow] (d) -- (e);
  \draw[arrow] (e) -- (f);
  \draw[arrow] (f) -- (g);
\end{tikzpicture}
\caption{Planning Agent pipeline.}
\end{figure}

\section{Weighted Scoring Algorithm}
\tacgia
\[
\operatorname{score}(m) =
w_d N_d(m) +
w_c N_c(m) +
w_w N_w(m) +
w_t N_t(m)
\]
Trong đó \(m\) là transport mode; \(N_d,N_c,N_w,N_t\) là duration, cost, walking và transfers đã chuẩn hóa. Trước khi tính tổng, weight được điều chỉnh theo mưa, giờ cao điểm và constraints.

\section{Adaptation and Failure Handling}
\tacgia
APScheduler chạy kiểm tra LTA mỗi hai phút và weather mỗi ba mươi phút. Khi cảnh báo ảnh hưởng trip, Adaptation Agent tạo proposal, tính delta và chờ người dùng chấp nhận. Một OneMap mode thất bại không làm hỏng toàn bộ request; Gemini và Supabase cũng có các nhánh fallback phù hợp.

\chapter{Implementation}
\section{Implementation Challenges and Solutions}
\tacgia
\begin{longtable}{p{3.5cm}p{5.5cm}p{6cm}}
\toprule
\textbf{Challenge} & \textbf{Solution} & \textbf{Remaining limitation} \\
\midrule
Multi-day scheduling & Day-bucketed greedy xét distance, dwell time và time window & Không bảo đảm tối ưu toàn cục \\
Multimodal routing & Fetch song song, normalize mode và fallback & Grab estimate không có surge realtime \\
Context-aware scoring & Relative normalization và weight adjustment & Chưa có crowding realtime \\
Live adaptation & Poll LTA/weather, tạo proposal và human approval & Scheduler chạy in-process \\
Place data quality & Startup validation và enrichment scripts & Cần tiếp tục audit dữ liệu/ảnh \\
\bottomrule
\end{longtable}

\chapter{Testing}
\section{Testing Strategy and Test Cases}
\tacgia
Dự án sử dụng Pytest cho backend và Vitest/Testing Library cho frontend. External APIs được mock trong phần lớn test.
\begin{longtable}{p{3.3cm}p{3cm}p{5cm}p{4cm}}
\toprule
\textbf{Feature} & \textbf{Level} & \textbf{Scenario} & \textbf{Expected result} \\
\midrule
Scoring & Unit & Duration-heavy profile & Fastest mode được ưu tiên \\
Scoring context & Unit & Heavy rain & Walking-heavy option giảm hạng \\
Day distribution & Unit & Multi-day và opening hours & Places phân bổ hợp lệ \\
OneMap & Service & Một mode lỗi & Mode khác vẫn trả về \\
Adaptation & Agent & MRT disruption hoặc rain & Tạo proposal \\
Planner UI & Component & Chọn preference & UI và request cập nhật đúng \\
\bottomrule
\end{longtable}

\section{Latest Test Results}
\tacgia
\begin{table}[H]
\centering
\begin{tabular}{lccc}
\toprule
\textbf{Suite} & \textbf{Passed} & \textbf{Failed} & \textbf{Date} \\
\midrule
Backend Pytest & Điền vào đây & Điền vào đây & Điền vào đây \\
Frontend Vitest & Điền vào đây & Điền vào đây & Điền vào đây \\
Frontend build & Điền vào đây & Điền vào đây & Điền vào đây \\
\bottomrule
\end{tabular}
\caption{Kết quả test gần thời điểm nộp bài.}
\end{table}
\danhinh{Ảnh kết quả test cuối cùng.}

\section{Bug Reports and Improvements}
\tacgia
Các vấn đề đã ghi nhận gồm Google OAuth provider chưa bật, localization tiếng Việt chưa đầy đủ, preference weights dạng phần trăm gây khó hiểu, map marker/control chưa rõ nghĩa và dữ liệu ảnh địa điểm cần tiếp tục audit.

\chapter{Demo}
\section{Demo Scenario}
\tacgia
\begin{enumerate}
  \item Mở Home page và tạo itinerary mới.
  \item Nhập ngày, khách sạn, ngân sách, preference và địa điểm.
  \item Tạo TripPlan và xem route được đề xuất.
  \item Xem Overview, Day view, map và alternatives.
  \item Đổi mode hoặc chỉnh thứ tự địa điểm.
  \item Kiểm tra cảnh báo và adaptation proposal.
  \item Gửi feedback và sử dụng chatbot.
\end{enumerate}

\section{Important UI Screenshots}
\tacgia
\danhinh{Trang Home/Landing của IMOVEV2.}
\danhinh{Planner: thiết lập itinerary và chọn địa điểm.}
\danhinh{Trip Overview và các metric chính.}
\danhinh{Day itinerary cùng bản đồ route.}
\danhinh{So sánh và đổi transport mode.}
\danhinh{Cảnh báo và adaptation proposal.}
\danhinh{Preference settings.}
\danhinh{Chatbot của IMOVEV2.}

\section{YouTube Demo Link}
\tacgia
\begin{tcolorbox}[colback=yellow!10,colframe=orange!70!black]
\centering\Large\textbf{Dán link video YouTube vào đây}
\end{tcolorbox}

\chapter{Deployment}
\section{Deployment Architecture}
\tacgia
Kiến trúc triển khai mục tiêu dùng Vercel cho frontend, Render cho FastAPI backend và Supabase cho PostgreSQL/Auth/Realtime.
\begin{figure}[H]
\centering
\begin{tikzpicture}[node distance=1.2cm and 1.2cm]
  \node[nodebox] (user) {User Browser};
  \node[external,right=of user] (vercel) {Vercel\\React Frontend};
  \node[external,right=of vercel] (render) {Render\\FastAPI Backend};
  \node[data,above right=of render] (supabase) {Supabase};
  \node[external,below right=of render] (apis) {OneMap / LTA /\\Weather / Gemini};
  \draw[arrow] (user) -- (vercel);
  \draw[arrow] (vercel) -- (render);
  \draw[arrow] (render) -- (supabase);
  \draw[arrow] (render) -- (apis);
\end{tikzpicture}
\caption{Kiến trúc deployment mục tiêu.}
\end{figure}

\section{Deployment Information}
\tacgia
\begin{itemize}
  \item Frontend URL: Điền vào đây.
  \item Backend URL: Điền vào đây.
  \item Environment variables và migration: Điền vào đây nếu có thay đổi.
\end{itemize}

\chapter{Logbook: Plan and Actual Workload}
\section{Planned and Actual Workload}
\tacgia
\begin{longtable}{p{2.7cm}p{3.2cm}p{5.5cm}p{2cm}p{2cm}}
\toprule
\textbf{Student ID} & \textbf{Full Name} & \textbf{Main Contributions} & \textbf{\%} & \textbf{Hours} \\
\midrule
Điền vào đây & Điền vào đây & Điền vào đây & Điền vào đây & Điền vào đây \\
Điền vào đây & Điền vào đây & Điền vào đây & Điền vào đây & Điền vào đây \\
Điền vào đây & Điền vào đây & Điền vào đây & Điền vào đây & Điền vào đây \\
Điền vào đây & Điền vào đây & Điền vào đây & Điền vào đây & Điền vào đây \\
\bottomrule
\end{longtable}
Tổng phần trăm actual workload phải bằng 100\%.

\section{Management Tools and Evidence}
\tacgia
Nhóm sử dụng Git/GitHub để quản lý source. Các công cụ giao tiếp, task management và tài liệu khác: Điền vào đây.
\danhinh{Bằng chứng công việc thành viên 1.}
\danhinh{Bằng chứng công việc thành viên 2.}
\danhinh{Bằng chứng công việc thành viên 3.}
\danhinh{Bằng chứng công việc thành viên 4.}

\chapter{AI usage and declaration}
\section{AI trong sản phẩm và quá trình phát triển}
\tacgia
IMOVE sử dụng Gemini 2.5 Flash cho parse tên địa điểm, đề xuất địa điểm, warning, gap notification và chatbot. Routing, cost, scheduling và scoring chính được thực hiện bằng code xác định và external routing data.
\begin{longtable}{p{3.2cm}p{4cm}p{4.5cm}p{4cm}}
\toprule
\textbf{AI Tool} & \textbf{Purpose} & \textbf{Output assisted} & \textbf{Human verification} \\
\midrule
Google Gemini & Product NLP & Parsing, suggestions, warning, chatbot & Validate structured output và fallback \\
Điền vào đây & Development assistance & Điền vào đây & Code review, tests và đối chiếu repository \\
\bottomrule
\end{longtable}

\section{Risk Control and Declaration}
\tacgia
Nhóm cam kết kiểm tra đầu ra AI trước khi sử dụng. Dữ liệu địa điểm do AI hỗ trợ có nguy cơ hallucination nên phải được audit bằng OneMap, Google Places hoặc nguồn đáng tin cậy. Các công cụ AI khác đã dùng và tuyên bố chính thức: Điền vào đây.

\chapter{Conclusion and Future Work}
\section{Conclusion}
\tacgia
IMOVEV2 đã xây dựng được web application lập itinerary nhiều ngày cho Singapore, kết hợp route đa phương thức, weighted scoring, map visualization, adaptation alerts, preference persistence và chatbot. Kiến trúc Agent/Service giúp phân tách trách nhiệm và giữ LLM ở vai trò hỗ trợ.

\section{Lessons Learned}
\tacgia
\begin{itemize}
  \item External data và API failure phải được xem là tình huống bình thường.
  \item Heuristic đơn giản nhưng đúng ràng buộc phù hợp với prototype.
  \item Typed contracts giảm lỗi tích hợp frontend/backend.
  \item AI hữu ích cho ngôn ngữ tự nhiên nhưng cần validation và fallback.
\end{itemize}

\section{Future Work}
\tacgia
\begin{itemize}
  \item Departure urgency engine và daily routine plan.
  \item Realtime crowding data trong scoring.
  \item Nâng chất lượng curated place data và ảnh.
  \item Provider abstraction cho nhiều AI model.
  \item Multi-city support.
  \item Hoàn thiện localization, preference UX, monitoring, CI/CD và user testing.
\end{itemize}

\section{Suggestions for Instructors and Teaching Assistants}
\tacgia
Điền vào đây.

\end{document}
